% element: 10-s066:e04
\BeleskeID{10-s066}
Sa prikazanim vrednostima komponenti tranzistor pri rastu izlaza dostiže 50\% ukupne promene dok je još u zasićenju. Zato kašnjenje treba povezati sa eksponencijalnim odzivom tog intervala. Dijagram posebno označava zasićeni i aktivni deo rasta.

Pri opadanju izlaz prazni približno stalna raspoloživa struja $I_{OL}$. Kondenzatorski odnos $i=C\,dv/dt$ tada daje linearan pad napona. Za polovinu promene važi
\[\frac{V_{OL}+V_{OH}}2=V_{OH}-\frac{I_{OL}}Ct_{pHL}.\]
Rešavanje po vremenu daje poslednji izraz. Njegova primena zahteva da struja u posmatranom intervalu ostane približno stalna.

Ekvivalent pražnjenja čine izlazna kapacitivnost $C_0$ i strujni izvor $\beta_FI_{B4}$, sa smerom pražnjenja prema masi. Ulazni impuls menja nivo od $V_H$ na $V_L$ u $t_1$, pa se vraća na $V_H$ u $t_2$.

U početku punjenja totem-pole radi u zasićenju, a zatim prelazi u aktivni režim. Pri pražnjenju prikazana zavisnost je linearna.
